% ============================================
% F 负责：详细设计 - 批阅模块 + Redis策略 + 部署
% ============================================

\section{批阅模块、缓存策略与部署详细设计}

本章对批阅子系统、Redis 缓存策略以及容器化部署方案进行详细设计。批阅子系统涵盖论文提交、教师批注、评语打分与退回重改四大核心流程；Redis 策略部分给出缓存 Key 规范、序列化方案及防穿透/击穿/雪崩的工程实现；部署部分提供完整的 Dockerfile、docker\-compose 编排文件与环境配置模板。

% ===== 7.1 批阅模块详细设计 =====
\subsection{批阅模块详细设计}

批阅模块是连接学生与教师的核心交互子系统。学生在完成论文编辑后提交论文进入批阅流程；教师对论文各章节添加批注，最终给出综合评语与打分（通过或退回）。本节从数据模型、服务层接口、RESTful API 与流程设计四方面展开。

\subsubsection{数据模型}

模块涉及三张核心数据表：\texttt{submission}（提交记录）、\texttt{annotation}（教师批注）和 \texttt{review\_record}（评语/批阅记录），以及关联的 \texttt{thesis}（论文）与 \texttt{thesis\_section}（章节）。

\begin{figure}[htbp]
\centering
\begin{verbatim}
thesis (论文)
   │  1
   │
   ├── N ── submission (提交记录)
   │           thesis_id  FK→thesis.id
   │           student_id FK→sys_user.id
   │           version_number
   │           submitted_at
   │
   ├── N ── annotation (批注)
   │           thesis_id  FK→thesis.id
   │           section_id FK→thesis_section.id
   │           teacher_id FK→sys_user.id
   │           start_offset, text_length
   │           selected_text, content
   │
   └── N ── review_record (批阅记录)
               thesis_id  FK→thesis.id
               teacher_id FK→sys_user.id
               comment_html, score, grade
               action: REVIEWED | RETURNED
               return_reason
\end{verbatim}
\caption{批阅模块实体关系图}
\label{fig:review-er}
\end{figure}

\textbf{（1）Submission —— 提交记录}

\begin{table}[htbp]
\centering
\caption{Submission 表字段说明}
\label{tab:submission-fields}
\begin{tabular}{|c|c|c|c|}
\hline
\textbf{字段} & \textbf{类型} & \textbf{约束} & \textbf{说明} \\
\hline
id & BIGSERIAL & PK & 主键自增 \\
\hline
thesis\_id & BIGINT & FK NOT NULL & 关联论文 \\
\hline
student\_id & BIGINT & FK NOT NULL & 提交学生 \\
\hline
version\_number & INT & NOT NULL, DEFAULT 1 & 提交版本号，每次提交递增 \\
\hline
submitted\_at & TIMESTAMP & DEFAULT NOW & 提交时间 \\
\hline
\end{tabular}
\end{table}

\textbf{（2）Annotation —— 教师批注}

\begin{table}[htbp]
\centering
\caption{Annotation 表字段说明}
\label{tab:annotation-fields}
\begin{tabular}{|c|c|c|c|}
\hline
\textbf{字段} & \textbf{类型} & \textbf{约束} & \textbf{说明} \\
\hline
id & BIGSERIAL & PK & 主键自增 \\
\hline
thesis\_id & BIGINT & FK NOT NULL & 关联论文 \\
\hline
section\_id & BIGINT & FK NOT NULL & 关联具体章节 \\
\hline
teacher\_id & BIGINT & FK NOT NULL & 批注教师 \\
\hline
start\_offset & INT & NOT NULL & 选中文本在章节内容中的起始偏移 \\
\hline
text\_length & INT & NOT NULL & 选中文本长度 \\
\hline
selected\_text & TEXT & 可空 & 被标注的原文内容 \\
\hline
content & TEXT & NOT NULL & 教师的批注意见 \\
\hline
created\_at & TIMESTAMP & DEFAULT NOW & 创建时间 \\
\hline
\end{tabular}
\end{table}

批注采用定位式设计：前端富文本编辑器在用户选中某段文字并添加批注时，记录选中文字相对于章节正文的 \texttt{start\_offset} 与 \texttt{text\_length}。重新渲染时根据偏移量定位到原文位置并高亮显示批注标记。

\textbf{（3）ReviewRecord —— 批阅记录}

\begin{table}[htbp]
\centering
\caption{ReviewRecord 表字段说明}
\label{tab:review-fields}
\begin{tabular}{|c|c|c|c|}
\hline
\textbf{字段} & \textbf{类型} & \textbf{约束} & \textbf{说明} \\
\hline
id & BIGSERIAL & PK & 主键自增 \\
\hline
thesis\_id & BIGINT & FK NOT NULL & 关联论文 \\
\hline
teacher\_id & BIGINT & FK NOT NULL & 批阅教师 \\
\hline
comment\_html & TEXT & 可空 & 富文本综合评语 \\
\hline
score & INT & CHECK [0,100] & 评分（0-100分制） \\
\hline
grade & VARCHAR(10) & 可空 & 等级映射（优/良/中/及格/不及格） \\
\hline
action & VARCHAR(20) & NOT NULL & REVIEWED（通过）或 RETURNED（退回） \\
\hline
return\_reason & TEXT & 可空 & 退回原因（action=RETURNED 时必填） \\
\hline
created\_at & TIMESTAMP & DEFAULT NOW & 创建时间 \\
\hline
\end{tabular}
\end{table}

\subsubsection{服务层设计}

为三层实体分别定义 Service 接口与实现类，统一继承项目已有的分层架构规范。

\textbf{SubmissionService 核心方法：}

\begin{lstlisting}[language=Java, caption=SubmissionService 接口核心方法]
public interface SubmissionService {
    /**
     * 学生提交论文：校验论文归属、更新状态为SUBMITTED、
     * 生成版本号递增的提交记录
     */
    Submission submitThesis(Long thesisId, Long studentId);

    /** 查询某论文的全部提交历史（按时间倒序） */
    List<Submission> getSubmissionHistory(Long thesisId);
}
\end{lstlisting}

提交逻辑实现要点：
\begin{enumerate}
    \item 校验 thesisId 对应的论文状态是否为 DRAFT 或 RETURNED（只有这两种状态允许提交）。
    \item 校验 studentId 是否为该论文的所有者（防止越权提交他人论文）。
    \item 查询当前最大 version\_number，新提交的版本号 = max + 1。
    \item 创建 Submission 记录，更新 thesis.status = SUBMITTED。
    \item 整个操作包裹在 \texttt{@Transactional} 中保证原子性。
\end{enumerate}

\textbf{AnnotationService 核心方法：}

\begin{lstlisting}[language=Java, caption=AnnotationService 接口核心方法]
public interface AnnotationService {
    /** 教师创建批注 */
    Annotation createAnnotation(Long thesisId, Long sectionId,
            Long teacherId, int startOffset, int textLength,
            String selectedText, String content);

    /** 教师修改批注内容 */
    Annotation updateAnnotation(Long annotationId, String newContent);

    /** 删除单个批注 */
    void deleteAnnotation(Long annotationId);

    /** 查询某论文全部批注（按创建时间排序） */
    List<Annotation> getAnnotationsByThesis(Long thesisId);

    /** 查询某论文某章节的全部批注（按偏移量排序） */
    List<Annotation> getAnnotationsBySection(Long thesisId, Long sectionId);
}
\end{lstlisting}

批注创建实现要点：
\begin{enumerate}
    \item 校验 teacherId 对应角色的用户是否为 TEACHER。
    \item 校验 sectionId 对应的章节是否属于 thesisId 对应的论文（防止跨论文越权批注）。
    \item start\_offset 与 text\_length 合法性校验：offset $\ge 0$，length $> 0$，offset + length $\le$ 章节正文总长度。
    \item 创建 Annotation 记录后返回给前端，前端根据偏移量在编辑器中渲染高亮标记。
\end{enumerate}

\textbf{ReviewRecordService 核心方法：}

\begin{lstlisting}[language=Java, caption=ReviewRecordService 接口核心方法]
public interface ReviewRecordService {
    /** 教师提交批阅结果 */
    ReviewRecord submitReview(Long thesisId, Long teacherId,
            String commentHtml, Integer score, String grade,
            String action, String returnReason);

    /** 查询某论文的全部批阅记录 */
    List<ReviewRecord> getReviewHistory(Long thesisId);
}
\end{lstlisting}

批阅逻辑实现要点：
\begin{enumerate}
    \item 校验论文状态必须为 SUBMITTED（已提交但尚未批阅的论文才可批阅）。
    \item 若 action = RETURNED，returnReason 不可为空，此时将 thesis.status 更新为 RETURNED。
    \item 若 action = REVIEWED，将 thesis.status 更新为 REVIEWED。
    \item score 范围校验：0 $\le$ score $\le$ 100。
    \item 整个操作在 \texttt{@Transactional} 中完成，确保评分记录与状态更新的一致。
\end{enumerate}

\subsubsection{API 端点设计}

所有 API 路径以 \texttt{/api/v1} 为前缀，遵循 RESTful 风格。需要认证的端点通过 \texttt{@RoleRequired} 注解声明所需角色。

\begin{table}[htbp]
\centering
\caption{批阅模块 API 端点一览}
\label{tab:review-api}
\begin{tabular}{|c|c|c|c|c|}
\hline
\textbf{方法} & \textbf{路径} & \textbf{角色} & \textbf{说明} \\
\hline
POST & /api/v1/submissions/{thesisId} & STUDENT & 提交论文 \\
\hline
GET & /api/v1/submissions/thesis/{thesisId} & ALL & 查看提交历史 \\
\hline
POST & /api/v1/annotations & TEACHER & 创建批注 \\
\hline
PUT & /api/v1/annotations/\{id\} & TEACHER & 修改批注 \\
\hline
DELETE & /api/v1/annotations/\{id\} & TEACHER & 删除批注 \\
\hline
GET & /api/v1/annotations/thesis/{thesisId} & ALL & 查看论文所有批注 \\
\hline
GET & /api/v1/annotations/thesis/{thesisId}/section/{sectionId} & ALL & 查看章节批注 \\
\hline
POST & /api/v1/reviews/{thesisId} & TEACHER & 提交批阅结果 \\
\hline
GET & /api/v1/reviews/thesis/{thesisId} & ALL & 查看批阅历史 \\
\hline
\end{tabular}
\end{table}

\subsubsection{批阅流程状态机}

论文在批阅子系统中经历以下状态迁移：

\begin{verbatim}
DRAFT ──(学生编辑完成)──→ COMPLETED ──(学生提交)──→ SUBMITTED
                                                       │
                    ┌──────────────────────────────────┤
                    │                                  │
                    ▼                                  ▼
              REVIEWED (通过)                    RETURNED (退回)
                                                      │
                                                      │
                    ┌─────────────────────────────────┘
                    │
                    ▼
           (学生修改论文，回到 DRAFT 状态，可再次提交)
\end{verbatim}

状态变更规则：
\begin{itemize}
    \item DRAFT → COMPLETED：学生确认论文编辑完成（前端触发，非提交动作）。
    \item COMPLETED → SUBMITTED：学生调用提交接口，创建 Submission 记录。
    \item SUBMITTED → REVIEWED：教师打回通过结果（action=REVIEWED），论文定稿。
    \item SUBMITTED → RETURNED：教师打回退回结果（action=RETURNED），须填写退回原因，学生可查看后修改并重新提交。
    \item RETURNED → SUBMITTED：学生修改后重新提交（视为下一版本）。
\end{itemize}

% ===== 7.2 Redis 缓存策略详细设计 =====
\subsection{Redis 缓存策略详细设计}

本节给出 Redis 在业务层中的具体应用模式与 Java 代码实现范例，涵盖缓存读取、写入、失效与防穿透/击穿/雪崩的工程手段。

\subsubsection{Cache-Aside 模式实现}

以学院列表缓存为例，展示标准的 Cache-Aside 读写流程：

\begin{lstlisting}[language=Java, caption=学院列表缓存的 Cache-Aside 实现]
@Service
@RequiredArgsConstructor
public class CollegeService {
    private final CollegeRepository collegeRepo;
    private final RedisTemplate<String, Object> redisTemplate;
    private static final String CACHE_KEY = "college:list";
    private static final Duration TTL = Duration.ofHours(1);

    public List<College> getAllColleges() {
        // 1. 查缓存
        Object cached = redisTemplate.opsForValue().get(CACHE_KEY);
        if (cached != null) {
            return (List<College>) cached;
        }
        // 2. 缓存未命中，查数据库
        List<College> colleges = collegeRepo.findAll();
        // 3. 写入缓存（带 TTL + 随机偏移防雪崩）
        if (!colleges.isEmpty()) {
            Duration ttl = TTL.plusMillis(
                ThreadLocalRandom.current().nextLong(0, 600_000)); // ±10min
            redisTemplate.opsForValue().set(CACHE_KEY, colleges, ttl);
        } else {
            // 4. 缓存空值防穿透，短 TTL
            redisTemplate.opsForValue().set(CACHE_KEY, "NULL",
                Duration.ofSeconds(60));
        }
        return colleges;
    }

    public College saveCollege(College college) {
        College saved = collegeRepo.save(college);
        // 5. 数据变更时主动删除缓存
        redisTemplate.delete(CACHE_KEY);
        return saved;
    }
}
\end{lstlisting}

\subsubsection{章节草稿自动保存}

学生编辑论文时，前端每隔 30 秒将当前章节的草稿内容通过 API 暂存至 Redis。Redis 的高写入性能天然适合此类高频小型写入场景，避免对 PostgreSQL 的频繁更新。

\begin{lstlisting}[language=Java, caption=草稿自动保存实现]
@Service
@RequiredArgsConstructor
public class DraftService {
    private final RedisTemplate<String, Object> redisTemplate;
    private static final Duration DRAFT_TTL = Duration.ofMinutes(30);

    public void saveDraft(Long thesisId, Long sectionId,
                          String draftContent) {
        String key = "draft:" + thesisId + ":" + sectionId;
        redisTemplate.opsForValue().set(key, draftContent, DRAFT_TTL);
    }

    public String getDraft(Long thesisId, Long sectionId) {
        String key = "draft:" + thesisId + ":" + sectionId;
        Object val = redisTemplate.opsForValue().get(key);
        return val != null ? val.toString() : null;
    }

    public void deleteDraft(Long thesisId, Long sectionId) {
        String key = "draft:" + thesisId + ":" + sectionId;
        redisTemplate.delete(key);
    }
}
\end{lstlisting}

\subsubsection{JWT 令牌管理}

用户登录成功后生成 JWT，将 Token 存储于 Redis 并与用户 ID 关联。用户主动登出或管理员强制下线时，将 Token 标记为失效。

\begin{lstlisting}[language=Java, caption=登录与登出的 Redis 操作]
// 登录成功
String token = jwtUtil.generateToken(userId, role);
redisTemplate.opsForValue().set("token:" + userId, token,
    Duration.ofHours(24));

// 登出
redisTemplate.delete("token:" + userId);
// Token 黑名单（可选增强：存储 jti 在黑名单中）
redisTemplate.opsForValue().set("blacklist:" + tokenId, "1",
    Duration.ofHours(24));
\end{lstlisting}

JwtFilter 在校验每个请求时，先从 Redis 查询 Token 是否在黑名单中。若存在，即使 JWT 签名有效也拒绝该请求。Token 续期则通过独立的 Refresh Token（有效期 7 天）进行，Refresh Token 同样存储于 Redis。

\subsubsection{API 限流实现}

基于 Redis 计数器实现登录接口的简单限流：

\begin{lstlisting}[language=Java, caption=基于 Redis 的登录限流实现]
@Component
@RequiredArgsConstructor
public class RateLimiter {
    private final RedisTemplate<String, Object> redisTemplate;

    // 60 秒内同一 IP 最多登录 5 次
    public boolean tryAcquire(String ip) {
        String key = "rate:" + ip + ":login";
        Long count = redisTemplate.opsForValue().increment(key);
        if (count != null && count == 1) {
            redisTemplate.expire(key, Duration.ofSeconds(60));
        }
        return count != null && count <= 5;
    }
}
\end{lstlisting}

每次登录尝试调用 \texttt{tryAcquire(clientIp)}，若返回 false 则拒绝请求并返回 HTTP 429（Too Many Requests）。

\subsubsection{缓存监控与运维}

\begin{itemize}
    \item \textbf{内存使用监控：} 通过 Redis \texttt{INFO memory} 命令定期采集 \texttt{used\_memory\_human} 指标，设置 80\% 告警阈值。
    \item \textbf{命中率统计：} 在 Service 层埋点记录缓存命中/未命中次数，通过 Actuator 的 \texttt{/actuator/metrics} 暴露自定义指标 \texttt{cache.hit.ratio}。
    \item \textbf{慢查询日志：} Redis 配置 \texttt{slowlog-log-slower-than 10000}（超过 10ms 记录），定期巡检有无因大 Key 或复杂数据结构导致的慢操作。
\end{itemize}

% ===== 7.3 部署详细设计 =====
\subsection{部署详细设计}

本节提供可直接使用的 Dockerfile、Docker Compose 编排文件及环境配置模板，确保在任何安装了 Docker 的 Linux 服务器上均可一键启动完整系统。

\subsubsection{Dockerfile}

采用多阶段构建（Multi-stage Build），第一阶段使用 Maven 编译打包，第二阶段使用精简 JRE 镜像运行，最终镜像体积控制在 200MB 以内。

\begin{lstlisting}[language=Dockerfile, caption=多阶段构建 Dockerfile]
# ============ 构建阶段 ============
FROM maven:3.9-eclipse-temurin-21 AS builder
WORKDIR /build
COPY pom.xml .
RUN mvn dependency:go-offline -B
COPY src ./src
RUN mvn clean package -DskipTests -q

# ============ 运行阶段 ============
FROM eclipse-temurin:21-jre-alpine
WORKDIR /app
COPY --from=builder /build/target/*.jar app.jar
# 健康检查
HEALTHCHECK --interval=30s --timeout=5s --retries=3 \
  CMD wget -qO- http://localhost:8080/actuator/health || exit 1
EXPOSE 8080
ENTRYPOINT ["java", "-jar", "app.jar"]
\end{lstlisting}

\subsubsection{Docker Compose 编排}

以下编排文件定义了三个服务及其依赖关系、网络与数据卷。

\begin{lstlisting}[language=YAML, caption=docker-compose.yml]
version: "3.9"
services:
  postgres:
    image: postgres:16-alpine
    container_name: thesis-pg
    environment:
      POSTGRES_DB: thesis_generator
      POSTGRES_USER: thesis_user
      POSTGRES_PASSWORD: ${DB_PASSWORD:-thesis_pass}
    volumes:
      - pgdata:/var/lib/postgresql/data
    networks:
      - thesis-net
    healthcheck:
      test: ["CMD-SHELL", "pg_isready -U thesis_user -d thesis_generator"]
      interval: 10s
      timeout: 5s
      retries: 5
    restart: unless-stopped
    deploy:
      resources:
        limits:
          cpus: "1.0"
          memory: 512M

  redis:
    image: redis:7-alpine
    container_name: thesis-redis
    command: redis-server --appendonly yes --maxmemory 256mb
      --maxmemory-policy allkeys-lru
    volumes:
      - redisdata:/data
    networks:
      - thesis-net
    healthcheck:
      test: ["CMD", "redis-cli", "ping"]
      interval: 10s
      timeout: 5s
      retries: 5
    restart: unless-stopped
    deploy:
      resources:
        limits:
          cpus: "0.5"
          memory: 256M

  app:
    build: .
    container_name: thesis-app
    ports:
      - "8080:8080"
    environment:
      SPRING_PROFILES_ACTIVE: prod
      SPRING_DATASOURCE_URL: jdbc:postgresql:
        //postgres:5432/thesis_generator
      SPRING_DATASOURCE_USERNAME: thesis_user
      SPRING_DATASOURCE_PASSWORD: ${DB_PASSWORD:-thesis_pass}
      SPRING_DATA_REDIS_HOST: redis
      SPRING_DATA_REDIS_PORT: 6379
      JWT_SECRET: ${JWT_SECRET}
    depends_on:
      postgres:
        condition: service_healthy
      redis:
        condition: service_healthy
    networks:
      - thesis-net
    restart: unless-stopped
    deploy:
      resources:
        limits:
          cpus: "2.0"
          memory: 1G

volumes:
  pgdata:
  redisdata:

networks:
  thesis-net:
    driver: bridge
\end{lstlisting}

\subsubsection{关键配置说明}

\begin{enumerate}
    \item \textbf{Redis 持久化：} \texttt{--appendonly yes} 启用 AOF（Append Only File）持久化，确保 Redis 重启后缓存数据不丢失（用于恢复草稿与令牌状态）。\texttt{--maxmemory 256mb --maxmemory-policy allkeys-lru} 设置内存上限并在满时自动淘汰最近最少使用的 Key。
    \item \textbf{数据库密码：} 通过 \texttt{\$\{DB\_PASSWORD\}} 环境变量注入，避免密码硬编码在配置文件中。启动时通过 \texttt{export DB\_PASSWORD=xxx \&\& docker-compose up -d} 或 \texttt{.env} 文件传入。
    \item \textbf{JWT 密钥：} 同样通过环境变量 \texttt{\$\{JWT\_SECRET\}} 注入，生产环境须更换为足够强度的随机字符串（建议 256 位以上）。
    \item \textbf{健康检查：} PostgreSQL 使用 \texttt{pg\_isready} 命令、Redis 使用 \texttt{redis-cli ping}、应用使用 Actuator health 端点。应用的 \texttt{depends\_on} 设置为依赖 PostgreSQL 与 Redis 均 healthy 后才启动。
    \item \textbf{网络隔离：} 仅应用容器的 8080 端口对外暴露，数据库与缓存的端口仅在内部 \texttt{thesis-net} 网络中可达，避免外部直接攻击数据存储层。
\end{enumerate}

\subsubsection{生产环境补充建议}

\begin{itemize}
    \item \textbf{反向代理：} 在 Docker Compose 中添加 Nginx 容器作为入口反向代理，配置 HTTPS 终止、静态资源缓存与请求限流。
    \item \textbf{日志收集：} 应用日志（通过 Logback）输出到 stdout，由 Docker 的 \texttt{json-file} 日志驱动收集，再通过 Filebeat 或 Promtail 接入集中式日志平台（ELK / Loki）。
    \item \textbf{监控告警：} 集成 Prometheus + Grafana，通过 Spring Boot Actuator 的 \texttt{/actuator/prometheus} 端点暴露 JVM 指标、HTTP 请求指标与自定义缓存命中率指标。
    \item \textbf{数据库备份：} 为 PostgreSQL 数据卷配置定时备份脚本（\texttt{pg\_dump}），将备份文件上传至对象存储（如阿里云 OSS / AWS S3）。
    \item \textbf{集群扩展：} 当单机性能不足时，可将应用容器扩展为多副本（\texttt{docker-compose up --scale app=3}），前端加 Nginx 负载均衡；Redis 升级为哨兵模式或集群模式；PostgreSQL 采用主从复制读写分离。
\end{itemize}

\subsubsection{部署操作手册}

\begin{enumerate}
    \item \textbf{环境准备：} 确保服务器已安装 Docker $\ge 24.0$ 与 Docker Compose $\ge 2.20$。
    \item \textbf{获取代码：} \texttt{git clone https://github.com/zzy154204-gif/thesis-generator.git \&\& cd thesis-generator}。
    \item \textbf{设置环境变量（可选）：}
    \begin{verbatim}
    export DB_PASSWORD=your_secure_db_password
    export JWT_SECRET=your_256bit_random_secret
    \end{verbatim}
    \item \textbf{构建并启动：} \texttt{docker-compose up -d --build}。
    \item \textbf{验证服务：}
    \begin{verbatim}
    curl http://localhost:8080/actuator/health
    # 预期返回: {"status":"UP"}
    \end{verbatim}
    \item \textbf{查看日志：} \texttt{docker-compose logs -f app}。
    \item \textbf{停止服务：} \texttt{docker-compose down}（保留数据）；\texttt{docker-compose down -v}（清除所有数据）。
\end{enumerate}
